% SPDX-License-Identifier: CC-BY-4.0
% Copyright (C) 2026 Fabian Januszewski, Abhijit (the Langlandscape team)
% GENERATED by tools/extract_statements.py from mathematics/chapters/ -- do not edit by hand.

\begin{quoted}{Proposition}{8.8}{79}{prop:s03-mazur}{Mazur check; seed Example mazur}{chapters/s03/1-graphs.tex}
For every prime $q$,
\[\Phicomb\big(\GDR_q(\hat\cO^\times)\big)\ \cong\ \Z\big/\mathrm{num}\tfrac{q-1}{12},\]
and this group is cyclic, generated via $\bar\lambda$ of Lemma~8.4. For $q\in\{2,3\}$ both sides are trivial.
\stpmc: Lemma~8.7, Prop.~8.3, Lemma~8.4. Lean (S03-L5, L6):
\lean{Langlandscape.Graph.mazurNum}, \lean{Langlandscape.Graph.mazur\_prod\_mul\_sum}, \lean{Langlandscape.Graph.mazur\_weights\_pairwise\_coprime}
\stformal{} (abstract arithmetic core: for weights in $\{1,2,3\}$ with the counts $n_2,n_3$ and $\sum w^{-1}=\frac{q-1}{12}$ one has
$\prod w\cdot\sum w^{-1}=\mathrm{num}\frac{q-1}{12}$, and the weights are pairwise coprime; $q\ge5$ only); the identification of the weights with the DR graph
(\lean{Langlandscape.Ribet.sum\_inv\_wx\_Ohat}, \lean{Langlandscape.Ribet.mazur\_check}, \stformal{} modulo the registered hypothesis \lean{Langlandscape.Hyp.EllipticCount} (and \pathtt{Hyp.EichlerMass})); the cases $q=2,3$ are \stformal{} as well, through the $q\le3$ clause of the same hypothesis \lean{Langlandscape.Hyp.EllipticCount} (which asserts $|\XB{\hat\cO^\times}|=1$ there; \lean{Langlandscape.Ribet.mazur\_check} covers every prime $q$), so the chapter's elementary derivation of that case is not formalized
\end{quoted}

\begin{quoted}{Proposition}{9.5}{83}{prop:s03-phieis}{$\Phicomb$ is Eisenstein}{chapters/s03/2-eisenstein.tex}
Let $q$ be any prime, $U=\Uzero M$, $B=B_q$, $\ell\nmid qM$, and let $T_\ell$ act on $\Z[\XBq U]$ by the Brandt matrix (homological
convention, $[x]\mapsto\sum_i[xg_i]$). Then $T_\ell$ is self-adjoint for $\langle,\rangle_w$, so it acts on $\Phicomb(\GDR_q(U))$, and
$T_\ell-(1+\ell)$ annihilates $\Phicomb(\GDR_q(U))$. \stproved{} (proof below, steps (1)--(4); valid for all primes $q$, including $q\in\{2,3\}$, and all
$M$ with $q\nmid M$ and $\ell\nmid qM$; it uses Lemma~9.2, Prop.~3.3(a),(b), Prop.~8.3, and the
finiteness of $\XBq U$, Prop.~1.6, which makes $\GDR_q(U)$ a finite graph). Lean: abstract core \stformal{}: \lean{Langlandscape.Graph.heckeQ} and \lean{Langlandscape.Graph.phiComb\_eisenstein} (steps (2)--(4), abstract: an integer matrix $B$ with row sums $\ell+1$,
$w_{e'}B(e,e')=w_eB(e',e)$ and $w_e\mid B(e,e')$ off the diagonal, acting on $\Phicomb$ of the two-vertex lattice $\Z[E]^0$ with weights $w_e>0$). Instantiation at $U=\hat\cO^\times$: \lean{Langlandscape.Ribet.phiComb\_eisenstein\_Ohat} \stformal{} modulo the registered hypothesis (step (1)) \lean{Langlandscape.Hyp.BrandtStabFree}.

Cross-references (not needed for the proof): [Ri90, Thm~3.12, p.~449; proof pp.~449--450, Remark~3.13] \stcited{}
\lit{ribet\_1990\_modular-representations-of-gal-qbar-q-arising-from-modular-forms}{Thm 3.12, p.~449}. Ribet's $\Phi=\coker(\iota)$ is our
$\Phicomb$, and his proof is the divisibility of the off-diagonal Brandt entries $B_\ell(i,j)$ by $e(i)=w_i$, which is step~(1) below, a purely quaternionic
fact. The geometric analogue for $\Phi_{N,p}$ is [Ed91, Th.~1, p.~163] \stcited{}
\lit{asterisque\_1991\_courbes-modulaires-et-courbes-de-shimura\_vol}{Edixhoven, Th\'eor\`eme~1, p.~163}, which assumes $p>3$.
\end{quoted}

\begin{quoted}{Theorem}{10.2}{89}{thm:s03-ribet}{definite Ribet sequence as a vertex collapse; seed Thm ribet}{chapters/s03/3-ribet.tex}
Let $q\nmid D$, $U$ as in \texttt{if:s02-quot}, $\G:=\GU\quot\cT_q$, and $|\G|$ its underlying graph, with the orientation (C1).
\begin{enumerate}[label=(\alph*)]
\item Over $\Z$ there is an exact sequence, equivariant for all Hecke operators $T=[UgU]$ with $g$ supported at a finite set $\Sigma\not\ni q$ of
  places at which $U$ is factorisable (Prop.~7.10(c); this holds for all standard levels):
  \[0\to H_1(|\G|,\Z)\to\Z[\XB{U'}]^0\xrightarrow{\ \alpha=(-s_{0*},\,s_{1*})\ }\big(\Z[\XB U]^0\big)^{2}\to0.\]
\item Let $p\ge5$ and let $V$ be a finitely generated $\Zp$-module with continuous $U_p$-action. Then, with the traces of (C3),
  \[0\to H_1(\GU,V)\to\Brandt{U'}V\xrightarrow{\alpha}\Brandt UV^{2}\to H_0(\GU,V)\to0.\]
\item $\rk_\Z H_1(|\G|)=|\XB{U'}|-2|\XB U|+1$. At a level where all $\Gamma_x$ are trivial, this is $(q-1)|\XB U|+1$, as in
  \pathtt{if:s02-indef}. For $V=\Zp$, $p\ge5$: $H_1(\GU,\Zp)=H_1(|\G|,\Z)\otimes\Zp$.
\end{enumerate}
Part (a) holds over every commutative ring $R$ in place of $\Z$ (with $H_1(|\G|,R)$).
\emph{Status.} Given (iii): (a), (c) first two sentences \stproved{} (Lemma~10.1 and Prop.~7.10); (b) and the last
sentence of (c) \stpmc{} (Prop.~7.18(b), which rests on the cited tree property and Shapiro's lemma; Lean \pathtt{Hyp.SerreCoh}). Without
assuming (iii): \stpmc{} via $\HypSA\wedge\HypHS\wedge\mathrm{Hyp.BTtree}\Rightarrow$(iii) (Prop.~7.10). The definite form of Ribet's
surjectivity is [Ri90, Thm~3.15, p.~451]
\lit{ribet\_1990\_modular-representations-of-gal-qbar-q-arising-from-modular-forms}{Thm 3.15, p.~451} \stcited. Ribet proves it with Eichler's
theorems. [Ri90, Rem.~3.18, p.~452]
\lit{ribet\_1990\_modular-representations-of-gal-qbar-q-arising-from-modular-forms}{Rem.~3.18, p.~452} records the graph-connectivity
proof (Livn\'e, Edixhoven), which is ours. Lean (S03-L9, any \texttt{CommRing R}, under \pathtt{Tree.BGConnected}, or under $\HypSA\wedge\HypHS\wedge\mathrm{Hyp.BTtree}$ in the \texttt{\_of\_SA\_HS} versions): \lean{Langlandscape.Ribet.ribet\_exact}, \lean{Langlandscape.Ribet.ribet\_rank}, \lean{Langlandscape.Ribet.ribet\_hecke\_equivariant}, \lean{Langlandscape.Ribet.boundary\_BG\_eq}, \lean{Langlandscape.Ribet.boundary\_BG\_apply}, \lean{Langlandscape.Ribet.ribet\_exact\_of\_SA\_HS}, \lean{Langlandscape.Ribet.ribet\_rank\_of\_SA\_HS} \stformal{} for (a) and the rank in (c);
coefficient layer for (b), definitions and compatibilities only (S03-L14): \lean{Langlandscape.Ribet.CoeffData}, \lean{Langlandscape.Ribet.CoeffData.ofLocal},
\lean{Langlandscape.Ribet.alphaV}, \lean{Langlandscape.Ribet.alphaVee}, \lean{Langlandscape.Ribet.s0TrV\_comp\_s0PullV}, \lean{Langlandscape.Ribet.alphaV\_triv},
\lean{Langlandscape.Ribet.alphaVee\_triv} \stformal{}; the exactness statement (b) with coefficients $V$ itself is not formalized (\leanplanned)
\end{quoted}

\begin{quoted}{Corollary}{10.5}{90}{cor:s03-ihara}{definite Ihara lemma, (i)--(ii); seed Cor.~ihara}{chapters/s03/3-ribet.tex}
(i) $\alpha$ in Thm~10.2(a) is surjective over $\Z$.
(ii) For every abelian group $A$ (in particular $A=\F_\ell$ for every prime $\ell$), the kernel of
$\alpha^\vee:(A^{\XB U})^2\to A^{\XB{U'}}$, $(f_0,f_1)\mapsto s_1^*f_1-s_0^*f_0$, is the line of constants $A\cdot(\mathbf1,\mathbf1)$.
(ii$'$) Let $p\ge5$ and let $W$ be a $\Zp$-module with continuous $U_p$-action. Then
\[\ker\big(\alpha^\vee:\Brandt UW^2\to\Brandt{U'}W\big)=\big\{(f_w,\,f_{[q]^{-1}w}):w\in W^{U^1_p}\big\},\qquad f_w(x):=[\tau(x)]^{-1}w,\]
and this is $H^0(\GU,W)=W^{U^1_p}$.
\stproved{} for (i), (ii) given (iii) (Lemma~10.1 and Prop.~7.10). (ii$'$): \stpmc{} (proof below, new: it checks
$f_w\in\Brandt UW$ by (C5) and $s_1^*f_{[q]^{-1}w}=s_0^*f_w$, and uses Lemma~7.17 (given (iii); the tree property is cited)
and the density of $\GU$ in $U^1_p$, i.e.\ Thm.~7.7; see (R4)). Here $U$ is assumed factorisable at $p$, $U=U_p\times U^{(p)}$ (true
for all standard levels), so that $\det:U_p\to\Zp^\times$ is onto and (C6) applies; ``continuous'' refers to a Hausdorff topology on $W$ for which
$W$ is a topological $\Zp$-module (e.g.\ $W$ finitely generated with the $p$-adic topology, or discrete). Lean (S03-L9): \lean{Langlandscape.Ribet.ihara\_surj}, \lean{Langlandscape.Ribet.ihara\_surj\_of\_SA\_HS} (i), \lean{Langlandscape.Ribet.ihara\_ker\_const}, \lean{Langlandscape.Ribet.ihara\_ker\_const\_of\_SA\_HS} (ii, any $A$) \stformal{}; (ii$'$) is not formalized.
\end{quoted}

\begin{quoted}{Theorem}{11.3}{95}{thm:s03-iharaH0}{Ihara (iii): $H_0$ and its Hecke action}{chapters/s03/4-ihara.tex}
Let $p\ge5$ and let $V$ be a finitely generated $\Zp$-module with continuous $U_p$-action. Use conventions (C1)--(C6).
\begin{enumerate}[label=(\alph*)]
\item The map
  \[\varepsilon:\Brandt UV^2\to V_{U^1_p},\qquad \varepsilon(f_0,f_1)=\sum_{[x]\in\XB U}w_x^{-1}\big([\tau(x)]\,f_0(x)+[q\,\tau(x)]\,f_1(x)\big)\]
  is well defined and surjective, and $\ker\varepsilon=\mathrm{im}\,\alpha$. Hence $\coker\alpha\cong H_0(\GU,V)=V_{U^1_p}$.
\item For $g$ supported at $\ell\nmid Dpq$, and $U$ factorisable at $\ell$, $U=U_\ell\times U^{(\ell)}$ (true for the standard levels,
  Def.~1.2), the induced action of $T_g$ on $\coker\alpha$ is
  \[T_g\ \longmapsto\ \deg(UgU)\cdot[\ell^{\,v_\ell(\Nrd g)}].\]
  In particular, $T_\ell$ acts as $(1+\ell)[\ell]$ and $S_\ell$ acts as $[\ell]^2$, the exponent being $+1$. If $U_p$ contains
  the scalars, then $[\ell]^2$ is the scalar $\ell$, which acts on $\Vk$ as $\ell^{k-2}$, in agreement with the classical
  normalisation $S_\ell=\ell^{k-2}$. The exponent $+1$ is for the cohomological operators $T_g$; for the transposed (homological) operators of the normative convention of stage 02, \pathtt{if:s02-indef}, it is $-1$, and there $S_\ell$ acts as $[\ell]^{-2}$ (Remark~11.4,
  exponent table, (R3)).
\item On the $\chi$-eigenpiece of $V_{U^1_p}\otimes\cOE$ for a character $\chi:\Zp^\times\to\cOE^\times$ (via $[\cdot]$), $T_\ell$ acts by
  $(1+\ell)\chi(\ell)$.
\end{enumerate}
\emph{Status.} \stpmc{} (proof below; the stage tex has \sttoprove{} with a sketch, which upgrades the seed's ``Hecke action \stsketch''). The
well-definedness of $\varepsilon$, $\varepsilon\circ\alpha=0$, surjectivity and the formula $\varepsilon\circ T_g=\deg(UgU)[\ell^m]\circ\varepsilon$ are \stproved{}
outright; $\ker\varepsilon\subseteq\mathrm{im}\,\alpha$ uses Cor.~10.5(ii$'$), hence (iii) and the density of $\GU$ in $U^1_p$ (full strong
approximation, (R4)); the identification with $H_0(\GU,V)$ is Prop.~7.18(b),(c). Lean (S03-L11): the definitions \lean{Langlandscape.Ribet.tau}, \lean{Langlandscape.Ribet.diamond}, \lean{Langlandscape.Ribet.epsilon} and the predicate
\lean{Langlandscape.Ribet.EpsWellDefined} \stformal{} (definitions only); clause (a) is the \emph{registered hypothesis} \lean{Langlandscape.Hyp.IharaH0} over
\pathtt{Ribet.alphaV} (an assumption, not a formalized theorem); clause (b) is
maths only here (the Hecke operators with coefficients, S04-L15a, now exist in Lean, but clause (b) itself is neither stated nor proved), \leanplanned
\end{quoted}

\begin{quoted}{Conjecture}{12.10}{107}{conj:s03-iharaGL2}{Ihara lemma for $\GL_2$, Betti form}{chapters/s03/5-amalgam.tex}
Let $\m$ be a maximal ideal of $\T^{(Mql)}$ (Lemma~12.7); call $\m$ \emph{Eisenstein} if there is $d\ge1$ such that $T_p-2\in\m$ for all but finitely
many primes $p\equiv1\bmod d$. This is the notion of [DT, p.~439--440]
\lit{diamond-taylor\_1994\_non-optimal-levels-of-mod-l-modular-representations}{p.~439--440}; DT's second condition $S_p-1\in\m$ is automatic for $p\equiv1\bmod dl$, since the
scalar $p$ acts as $p^{k-2}$ (Lemma~12.4(b)) and $l\in\m$. Compare Cor.~11.5 on the quaternionic side.
For non-Eisenstein $\m$, $(\coker\Psi)_{\m,\mathrm{tors}}=0$. Equivalently, $\Psi\otimes\F_l$ is injective after localisation at
$\m$, because the groups $H^1(\cdot,V)$ are torsion-free (so the Tor sequence of $\otimes\F_l$ has no extra term; $H^2(\Gamma,V/l)$ plays no role) and $\Psi$ is injective (Lemma~12.8). \stexp{} (consistent with C3: all torsion found is
Eisenstein, at $l=k+1$). The classical route uses Ihara's lemma [Ih75] (not in \pathtt{literature/}; \texttt{ihara\_1966} there is a different paper) or the congruence subgroup
property of $\SL_2(\Z[\frac1q])$ (\pathtt{literature/serre\_1970\_*}, scan, theorem not checked). The indefinite quaternionic analogue
is [MS19, Thm~1.1]
\lit{manning-shotton\_2019\_iharas-lemma-for-shimura-curves-over-totally-real-fields-via-patching\_arxiv}{Thm 1.1, p.~2} \stcited.
The ``equivalently'' needs $\ker\Psi=H^1(G,V)=0$, which is Lemma~12.8. Writing $0\to K\to A\xrightarrow\Psi B\to C\to0$ with
$A=H^1(\Gamma,V)^2$, $B=H^1(\Gamma_0(Mq),V)$ free over $\Z_{(l)}$ (Prop.~12.3), one has $\ker(\Psi\otimes\F_l)\supseteq K/l$ and
$\ker(\Psi\otimes\F_l)/(K/l)\cong C[l]$; so, as $K=0$, $\Psi\otimes\F_l$ is injective after localisation at $\m$ iff $C[l]_\m=0$ iff $C_{\m,\mathrm{tors}}=0$
(all maps are $\T^{(Mql)}$-linear by Lemma~12.7).
\end{quoted}
